\documentclass[11pt]{beamer}
\usetheme{CambridgeUS}
% =====================
% Packages
% =====================
\usepackage[french]{babel}
\usepackage[T1]{fontenc}
\usepackage{lmodern}
\usepackage{amsmath,amssymb}
\usepackage{graphicx}
% =====================
% Titre / Auteurs
% =====================
\title{Modélisation de la sévérité et de la fréquence des accidents routiers}
\author[Equipe 15]{%
BKHACH Ayoub\\
BEGGAR Rayane\\
SABRI Yasser\\
MARSOUK Ayoub
}
\institute[ENSIMAG]{%
ENSIMAG\\[0.5em]
\textbf{Encadrant :} Christophe Dutang
}
\date{\today}
% =====================
% Début document
% =====================
\begin{document}
% ---------- Page de titre ----------
\begin{frame}
\titlepage
\end{frame}

% ---------- Déclaration de TOUTES les sections d'abord ----------
\section{Introduction}
\section{Modèle linéaire et score de sévérité}
\section{Les regroupements avec la méthode de partitionnement}
\section{Conclusion}

% ---------- Plan ----------
\begin{frame}{Plan}
\tableofcontents
\end{frame}

% =====================
% Introduction
% =====================
\subsection*{} % Pour revenir à la section Introduction
\begin{frame}{Contexte}
\begin{itemize}
  \item Problématique des accidents routiers
  \item Importance de la modélisation statistique
  \item Objectifs de l'étude
\end{itemize}
\end{frame}

% =====================
% Modèle linéaire et score
% =====================
\subsection*{} % Pour passer à la section suivante
\begin{frame}{Score de sévérité}
\begin{itemize}
  \item Définition d'un score de sévérité
  \item Combinaison des blessés, hospitalisations et décès
\end{itemize}
\end{frame}

\begin{frame}{Modèle linéaire}
\begin{itemize}
  \item Variable réponse : score de sévérité
  \item Variables explicatives : contexte, infrastructure, véhicule
\end{itemize}
\end{frame}

\begin{frame}{Modèle linéaire simplifié}
\begin{itemize}
  \item Sélection des variables significatives
  \item Amélioration de l'interprétabilité
\end{itemize}
\end{frame}

% =====================
% Regroupements
% =====================
\subsection*{} % Pour passer à la section suivante
\begin{frame}{Méthode de partitionnement}
\begin{itemize}
  \item Clustering des accidents
  \item Mise en évidence de profils types
\end{itemize}
\end{frame}

% =====================
% Conclusion
% =====================
\subsection*{} % Pour passer à la section Conclusion
\begin{frame}{Conclusion}
\begin{itemize}
  \item Résultats principaux
  \item Limites de l'étude
  \item Perspectives
\end{itemize}
\end{frame}

\end{document}